\SourceLeaf{252}%
\(\mathrm{d}p,\mathrm{d}q,\mathrm{d}p',\mathrm{d}q'\), respondentes erunt
\[
\BareRoot{E}\mathord{.}\mathrm{d}p,\qquad\BareRoot{G}\mathord{.}\mathrm{d}q,\qquad\BareRoot{E'}\mathord{.}\mathrm{d}q',\qquad\BareRoot{G'}\mathord{.}\mathrm{d}q'
\]
Angulos, quos horum elementorum directiones faciunt cum directione fixa arbitraria, denotabimus per \(M,N,M',N'\), numerando eo sensu, quo iacet secunda respectu primae, ita ut \(\sin(N-M)\) fiat quantitas positiva: eodem sensu iacere supponemus (quod licet) quartam respectu tertiae, ita ut etiam \(\sin(N'-M')\) sit quantitas positiva. His ita intellectis, si consideramus punctum aliud, a priore infinite parvum distans, cui respondeant valores variabilium
\[
p+\mathrm{d}p,\quad q+\mathrm{d}q,\quad p'+\mathrm{d}p',\quad q'+\mathrm{d}q'
\]
levi attentione adhibita cognoscemus, fieri generaliter, i. e. independenter a valoribus variationum \(\mathrm{d}p,\mathrm{d}q,\mathrm{d}p',\mathrm{d}q'\),
\[
\begin{aligned}
&\BareRoot{E}\mathord{.}\mathrm{d}p\mathord{.}\sin M+\BareRoot{G}\mathord{.}\mathrm{d}q\mathord{.}\sin N\\
&=\BareRoot{E'}\mathord{.}\mathrm{d}p'\mathord{.}\sin M'+\BareRoot{G'}\mathord{.}\mathrm{d}q'\mathord{.}\sin N'
\end{aligned}
\]
quum utraque expressio nihil aliud sit, nisi distantia puncti novi a linea, a qua anguli directionum incipiunt. Sed habemus, per notationem iam supra introductam \(N-M=\omega\). et per analogiam statuemus \(N'-M'=\omega'\), nec non insuper \(N-M'=\psi\). Ita aequatio modo inventa exhiberi potest in forma sequenti
\[
\begin{aligned}
&\BareRoot{E}\mathord{.}\mathrm{d}p\mathord{.}\sin(M'-\omega+\psi)+\BareRoot{G}\mathord{.}\mathrm{d}q\mathord{.}\sin(M'+\psi)\\
&=\BareRoot{E'}\mathord{.}\mathrm{d}p'\mathord{.}\sin M'+\BareRoot{G'}\mathord{.}\mathrm{d}q'\mathord{.}\sin(M'+\omega')
\end{aligned}
\]
vel ita
\[
\begin{aligned}
&\BareRoot{E}\mathord{.}\mathrm{d}p\mathord{.}\sin(N'-\omega-\omega'+\psi)+\BareRoot{G}\mathord{.}\mathrm{d}q\mathord{.}\sin(N'-\omega'+\psi)\\
&=\BareRoot{E'}\mathord{.}\mathrm{d}p'\mathord{.}\sin(N'-\omega')+\BareRoot{G'}\mathord{.}\mathrm{d}q'\mathord{.}\sin N'
\end{aligned}
\]

Et quum aequatio manifesto independens esse debeat a directione initiali, hanc ad lubitum accipere licet. Statuendo itaque in forma secunda \(N'=0\), vel in prima \(M'=0\), obtinemus aequationes sequentes:
\[
\begin{aligned}
\BareRoot{E'}\mathord{.}\sin\omega'\mathord{.}\mathrm{d}p'&=\BareRoot{E}\mathord{.}\sin(\omega+\omega'-\psi)\mathord{.}\mathrm{d}p+\BareRoot{G}\mathord{.}\sin(\omega'-\psi)\mathord{.}\mathrm{d}q\\
\BareRoot{G'}\mathord{.}\sin\omega'\mathord{.}\mathrm{d}q'&=\BareRoot{E}\mathord{.}\sin(\psi-\omega)\mathord{.}\mathrm{d}p+\BareRoot{G}\mathord{.}\sin\psi\mathord{.}\mathrm{d}q
\end{aligned}
\]
quae aequationes quum identicae esse debeant cum his
\[
\begin{aligned}
\mathrm{d}p'&=\alpha\mathrm{d}p+\beta\mathrm{d}q\\
\mathrm{d}q'&=\gamma\mathrm{d}p+\delta\mathrm{d}q
\end{aligned}
\]
